\subsection{Dinamika Pelaksanaan UUD NRI Tahun 1945}

\subsubsection{Masa Awal Kemerdekaan dan Era Kabinet}
\begin{itemize}
\item UUD NRI Tahun 1945 pada awal kemerdekaan belum dapat dilaksanakan dengan baik karena negara terfokus pada perang mempertahankan kemerdekaan dari Belanda serta mengatasi pemberontakan internal, seperti PKI Madiun (18 September 1948) dan DI/TII Jawa Barat (7 Agustus 1949).
\item Penerbitan Maklumat 14 November 1945 mengubah sistem pemerintahan presidensial menjadi sistem parlementer yang dipimpin oleh Perdana Menteri.
\item Terdapat beberapa kali pergantian kabinet dalam rentang waktu 14 November 1945 hingga 27 Desember 1949:
  \begin{enumerate}
  \item Kabinet Sjahrir (14 November 1945--27 Juni 1947, dengan periode 27 Juni--3 Juli 1947 kembali sementara ke Kabinet Presidensial).
  \item Kabinet Amir Sjarifuddin I (3 Juli--11 November 1947).
  \item Kabinet Amir Sjarifuddin II (11 November 1947--29 Januari 1948).
  \item Kabinet Mohammad Hatta I (29 Januari 1948--4 Agustus 1949).
  \item Kabinet Syafruddin Prawiranegara / Kabinet Darurat (19 Desember 1948--13 Juli 1949).
  \item Kabinet Mohammad Hatta II (4 Agustus 1949--20 Desember 1949).
  \end{enumerate}
\end{itemize}

\subsubsection{Masa Konstitusi RIS}
\begin{itemize}
\item Sebagai konsekuensi Konferensi Meja Bundar (KMB), UUD NRI Tahun 1945 digantikan oleh Konstitusi Republik Indonesia Serikat (RIS) pada tanggal 27 Desember 1949.
\item Bentuk negara berubah dari negara kesatuan menjadi negara serikat (federasi), dengan perubahan pada struktur kelembagaan dan hubungan pusat-daerah.
\item Konstitusi RIS berlaku di seluruh wilayah federasi, sementara UUD NRI Tahun 1945 tetap berlaku khusus bagi wilayah negara bagian Republik Indonesia (beribu kota di Yogyakarta).
\item Tuntutan rakyat untuk kembali ke bentuk kesatuan menyebabkan berlakunya Konstitusi RIS hanya bertahan sampai 17 Agustus 1950.
\end{itemize}

\subsubsection{Masa UUDS 1950}
\begin{itemize}
\item Berlakunya Undang-Undang Dasar Sementara (UUDS) 1950 mengembalikan bentuk negara menjadi negara kesatuan dengan sistem kabinet parlementer.
\item Masa ini ditandai dengan berkembangnya kehidupan kepartaian dan politik multipartai.
\item Sistem ini berakhir setelah dikeluarkannya Dekrit Presiden pada 5 Juli 1959.
\end{itemize}

\subsubsection{Masa Demokrasi Terpimpin dan Orde Lama}
\begin{itemize}
\item Dekrit Presiden 5 Juli 1959 memuat tiga ketentuan utama:
  \begin{enumerate}
  \item Pembubaran Konstituante.
  \item Pemberlakuan kembali UUD NRI Tahun 1945 dan tidak berlakunya UUDS 1950.
  \item Pembentukan MPRS dan DPAS dalam waktu sesingkat-singkatnya.
  \end{enumerate}
\item Pelaksanaan UUD NRI Tahun 1945 pada era ini mengalami beberapa penyimpangan, antara lain:
  \begin{itemize}
  \item Penerbitan Penetapan Presiden (Penpres) yang kedudukannya disetarakan dengan undang-undang tanpa persetujuan DPR, seperti Penpres No. 2/1959 (pembentukan MPRS), Penpres No. 7/1959 (pembubaran partai politik), Penpres No. 1/1960 (Manipol sebagai GBHN), Penpres No. 3/1960 (pembubaran DPR hasil Pemilu 1955), dan Penpres No. 4/1960 (pembentukan DPR-GR).
  \item Keputusan MPRS yang menyimpang dari konstitusi, seperti Tap MPRS No. III/MPRS/1963 yang mengangkat Presiden Sukarno sebagai Presiden Seumur Hidup.
  \item Pemusatan kekuasaan pada Presiden sehingga menempatkan DPR di bawah kedudukan Presiden.
  \end{itemize}
\end{itemize}

\subsubsection{Masa Orde Baru}
\begin{itemize}
\item Dimulai sejak diterbitkannya Surat Perintah 11 Maret 1966 (Supersemar) dengan ikhtiar melaksanakan Pancasila dan UUD NRI Tahun 1945 secara murni dan konsekuen.
\item Terjadi beberapa penyimpangan konstitusional selama masa Orde Baru:
  \begin{itemize}
  \item Perekonomian terpusat dan berpraktik monopoli yang kurang sesuai dengan asas Pasal 33 UUD 1945.
  \item Kekuasaan Presiden sangat dominan melebihi lembaga legislatif (executive heavy).
  \item Pengekangan kebebasan pers (pembredelan media) serta pembatasan hak berpendapat dan berorganisasi.
  \item Maraknya praktik Korupsi, Kolusi, dan Nepotisme (KKN).
  \item Penegakan hukum yang belum berjalan baik dan rentan diintervensi oleh kekuasaan.
  \end{itemize}
\end{itemize}

\subsubsection{Masa Reformasi}
\begin{itemize}
\item Bergulir sejak 21 Mei 1998 untuk menata ulang sistem ketatanegaraan dan mencegah pemusatan kekuasaan.
\item Alasan mendasar dilakukannya amendemen UUD NRI Tahun 1945:
  \begin{itemize}
  \item Pasal-pasal UUD 1945 sebelum amendemen bersifat multitafsir dan memberi peluang penafsiran ganda bagi pemegang kekuasaan.
  \item Pertanggungjawaban Presiden kepada MPR membuat Presiden secara tidak langsung tidak bertanggung jawab kepada DPR.
  \item Perlunya memberikan otonomi daerah yang luas kepada pemerintah daerah.
  \item Perlunya wadah hukum yang lebih luas dan tegas bagi perlindungan Hak Asasi Manusia (HAM).
  \end{itemize}
\end{itemize}

\subsection{Pergantian dan Perubahan Konstitusi}

\subsubsection{Sifat dan Prosedur Perubahan Konstitusi}
\begin{itemize}
\item Konstitusi memiliki dua sifat utama:
  \begin{itemize}
  \item Rigid (kaku): Memerlukan syarat dan tata cara khusus untuk diubah. UUD NRI Tahun 1945 bersifat rigid (ketentuan perubahan diatur dalam Pasal 37).
  \item Fleksibel (luwes): Dapat diubah dengan prosedur yang sama seperti pembuatan undang-undang biasa.
  \end{itemize}
\item Pembatasan perubahan UUD NRI Tahun 1945 diatur dalam Pasal 37 Ayat (5), yang menetapkan bahwa bentuk Negara Kesatuan Republik Indonesia tidak dapat diubah.
\item Cara perubahan konstitusi di berbagai negara secara umum dapat dilakukan melalui referendum, keputusan badan legislatif dengan syarat khusus, lembaga khusus pembentuk konstitusi, atau persetujuan negara-negara bagian dalam negara federasi.
\end{itemize}

\subsubsection{Kesepakatan Dasar Perubahan UUD NRI Tahun 1945}
\begin{enumerate}
\item Tidak mengubah Pembukaan UUD NRI Tahun 1945.
\item Tetap mempertahankan Negara Kesatuan Republik Indonesia (NKRI).
\item Mempertegas sistem pemerintahan presidensial.
\item Memindahkan hal-hal normatif dalam Penjelasan UUD NRI Tahun 1945 ke dalam pasal-pasal.
\item Melakukan perubahan dengan cara addendum.
\end{enumerate}

\subsubsection{Perbandingan Sebelum dan Sesudah Amendemen UUD NRI Tahun 1945}
\begin{itemize}
\item Sistematika: Sebelum amendemen terdiri dari 16 bab, 37 pasal, 49 ayat, 4 pasal Aturan Peralihan, dan 2 ayat Aturan Tambahan. Setelah amendemen menjadi 21 bab, 73 pasal, 170 ayat, 3 pasal Aturan Peralihan, dan 2 ayat Aturan Tambahan.
\item Kedaulatan Rakyat: Sebelum amendemen dilakukan sepenuhnya oleh MPR. Setelah amendemen dilaksanakan menurut UUD.
\item Kekuasaan Presiden: Berubah dari kekuasaan yang terpusat (executive heavy) menjadi penerapan prinsip saling mengawasi dan mengimbangi (check and balances).
\item Pemilihan Presiden: Sebelum amendemen dipilih oleh MPR, setelah amendemen dipilih langsung oleh rakyat melalui Pemilihan Umum.
\item Kelembagaan Negara: Kedudukan lembaga negara menjadi sejajar di bawah UUD. Dewan Pertimbangan Agung (DPA) dibubarkan, serta dibentuk lembaga baru yaitu Mahkamah Konstitusi (MK), Komisi Yudisial (KY), dan Dewan Perwakilan Daerah (DPD).
\item Keanggotaan MPR: Terdiri atas anggota DPR dan anggota DPD yang seluruhnya dipilih melalui pemilu.
\item Hak Asasi Manusia: Penjabaran HAM diatur lebih terperinci dan tegas dalam Pasal 28A sampai Pasal 28J.
\end{itemize}

\subsection{Peran Mahkamah Konstitusi}

\subsubsection{Sejarah Pembentukan dan Status}
\begin{itemize}
\item Pembentukan Mahkamah Konstitusi (MK) diawali dari pengadopsian ide MK dalam amendemen UUD tahun 2001 (Pasal 24 ayat 2, Pasal 24C, dan Pasal 7B).
\item Landasan operasional diatur melalui Undang-Undang Nomor 24 Tahun 2003 tentang Mahkamah Konstitusi yang disahkan pada 13 Agustus 2003.
\item Peradilan MK merupakan pengadilan tingkat pertama dan terakhir yang putusannya bersifat final.
\end{itemize}

\subsubsection{Kewenangan dan Kewajiban MK}
\begin{itemize}
\item Kewenangan MK berdasarkan Pasal 24C Ayat (1):
  \begin{enumerate}
  \item Menguji undang-undang terhadap UUD NRI Tahun 1945 (judicial review).
  \item Memutus sengketa kewenangan lembaga negara yang kewenangannya diberikan oleh UUD.
  \item Memutus pembubaran partai politik.
  \item Memutus perselisihan tentang hasil pemilihan umum.
  \end{enumerate}
\item Kewajiban MK: Wajib memberikan putusan atas pendapat DPR mengenai dugaan pelanggaran hukum atau tindak pidana oleh Presiden dan/atau Wakil Presiden.
\end{itemize}

\subsubsection{Peran Utama MK}
\begin{enumerate}
\item Pengawal konstitusi (the guardian of constitution).
\item Penafsir akhir konstitusi (the final interpreter of constitution).
\item Pengawal demokrasi (the guardian of democracy).
\item Pelindung hak-hak konstitusional warga negara (the protector of citizen's constitutional rights).
\item Pelindung hak asasi manusia (the protector of human rights).
\end{enumerate}

\subsection{Hak dan Kewajiban Warga Negara}

\subsubsection{Pengertian dan Sifat Hubungan Hak dan Kewajiban}
\begin{itemize}
\item Warga Negara Indonesia adalah orang-orang bangsa Indonesia asli dan orang-orang bangsa lain yang disahkan dengan undang-undang sebagai warga negara.
\item Sifat hubungan hak dan kewajiban warga negara:
  \begin{itemize}
  \item Timbal balik dengan negara: Hak warga negara merupakan kewajiban negara untuk memenuhinya, sedangkan kewajiban warga negara merupakan hak negara untuk menuntutnya.
  \item Kausalitas (sebab-akibat): Hak seseorang diperoleh sebagai akibat dari terlaksananya kewajiban.
  \item Tidak dapat dipisahkan: Hak dan kewajiban melekat secara beriringan.
  \end{itemize}
\end{itemize}

\subsubsection{Pasal-Pasal Hak dan Kewajiban dalam UUD NRI Tahun 1945}
\begin{itemize}
\item Pasal 27 Ayat (1): Kesamaan kedudukan warga negara di dalam hukum dan pemerintahan serta kewajiban menjunjung hukum dan pemerintahan.
\item Pasal 27 Ayat (2): Hak atas pekerjaan dan penghidupan yang layak bagi kemanusiaan.
\item Pasal 28: Kemerdekaan berserikat, berkumpul, serta mengeluarkan pikiran dengan lisan dan tulisan.
\item Pasal 28A: Hak untuk hidup serta mempertahankan hidup dan kehidupan.
\item Pasal 28B: Hak membentuk keluarga melalui perkawinan sah, melanjutkan keturunan, serta hak anak atas kelangsungan hidup, tumbuh berkembang, dan perlindungan dari kekerasan/diskriminasi.
\item Pasal 28C: Hak mengembangkan diri, memperoleh pendidikan, memanfaatkan ilmu pengetahuan dan teknologi, seni dan budaya, serta memajukan diri secara kolektif.
\item Pasal 28D: Hak atas pengakuan dan kepastian hukum yang adil, perlakuan sama di hadapan hukum, imbalan dan perlakuan adil dalam hubungan kerja, kesempatan sama dalam pemerintahan, serta status kewarganegaraan.
\item Pasal 28E: Kebebasan memeluk agama, memilih pendidikan, pekerjaan, kewarganegaraan, tempat tinggal, serta kebebasan meyakini kepercayaan dan menyatakan pikiran/sikap.
\item Pasal 28I: Jaminan atas hak-hak yang tidak dapat dikurangi dalam keadaan apa pun (non-derogable rights), seperti hak untuk hidup, tidak disiksa, tidak diperbudak, dan tidak dituntut berdasarkan hukum yang berlaku surut.
\item Pasal 28J: Kewajiban setiap orang untuk menghormati hak asasi manusia orang lain dan tunduk pada pembatasan undang-undang demi menjamin pengakuan serta penghormatan atas hak orang lain.
\end{itemize}

\subsubsection{Jaminan Kebebasan Beragama dan Berkepercayaan}
\begin{itemize}
\item Landasan konstitusional kebebasan beragama bersumber dari Pembukaan UUD NRI Tahun 1945 alinea ketiga dan keempat (Ketuhanan Yang Maha Esa) serta Pasal 29 UUD NRI Tahun 1945.
\item Negara menjamin kemerdekaan tiap-tiap penduduk untuk memeluk agamanya masing-masing dan untuk beribadat menurut agama dan kepercayaannya itu.
\end{itemize}
